% 7.模型评价.tex —— 七、模型的评价
\section{模型的评价}

\subsection{模型的优点}
\begin{itemize}
  \item \textbf{四问关系明确：}问题一的解耦模型用于预热期分析，问题二通过状态相关物性描述热质耦合，问题三在同一全过程模型上增加全场达标判据，问题四在附录 4 物性之上进一步引入附件 2 的移动边界，逐问扩展物理机制而避免重复建立模型。
  \item \textbf{离散结构可复核：}中心、内部与表面均从径向通量平衡构造，内部界面通量相消。问题一的解析算例、网格收敛与离散平衡检查分别支持实现正确性、精度和守恒结构；这些检查的适用范围可以明确追溯。
  \item \textbf{效率与误差分别考察：}显式 Euler 提供基线，CN 与三对角求解降低细空间网格下的时间步限制；问题三通过时间推进、输出和空间网格的分项检查，识别了所测设置下终止时间主要受空间分辨率影响；问题四的空间与时间收敛均达秒级一致，且收敛明显快于问题三，说明移动边界模型的归一化网格不存在收缩末期的边界层跨网格问题。
\end{itemize}

\subsection{模型的局限}
\begin{itemize}
  \item \textbf{能量与迁移模型仍有简化：}未计汽化潜热、迁移焓与内部蒸汽迁移，后两问也未完整追踪组分变化。遗漏蒸发耗热可能使温度预测偏高，但现有证据不足以定量确定其偏差，也不能据此认定它必然是最大误差源。
  \item \textbf{空间适用范围有限：}问题一至三的固定径向模型忽略收缩，问题四虽引入径向收缩，但把附件 2 的 $R(t)$ 作为已知边界直接使用，未建立“水分流失量 $\to$ 收缩量”的本构关系，且轴向收缩与端部传递仍未建模。短时中段近似的理由不能直接保证数天干燥过程中整个有限长药材都满足相同精度。
  \item \textbf{数值结论具有条件：}CN 的无条件稳定性限于相应线性系统，不能替代非线性残差、正性和大步长精度检查。含 Newton 敏度项的块交替迭代不保证整体二次收敛；子步数或迭代次数达到上限时，设定容差也未必得到满足。
  \item \textbf{实物预测尚待验证：}经验物性、传质边界和 $4$ h 后环境末值延拓会影响预测。目前的网格与步长检验说明数值解的自洽程度，不能代替实测温度、含水率和干燥时间的验证。约 $57.17$ h 是当前假设和参数下的趋势估计，而 $3431$ min 是其按分钟向上取整得到的保守口径；分钟化建议不意味着实际过程具备分钟级精度。
  \item \textbf{验证范围有限：}解析解只验证了简化线性导热算例的数值实现，不能作为完整干燥模型的现实精度证明；在忽略汽化潜热、迁移焓与内部蒸汽迁移的前提下，结果更适合作为相对比较与工艺时间估计。
  \item \textbf{QS 证据并不独立：}形状漂移参数使用起点后的同一数值轨道拟合，因此 QS 比较属于事后降阶一致性检查，不能用作独立预测或独立校核终止时间。
\end{itemize}
